\section{Methodology}
\label{sec:methodology}

The AURA architecture decouples perceptual state representation, policy optimization, and action-duration bounding into modular stages, as illustrated in Fig.~\ref{fig:architecture}.

\begin{figure}[htbp]
\centering
\begin{tikzpicture}[
    node distance=0.7cm and 0.5cm,
    block/.style={rectangle, draw=blue!80, fill=blue!5, thick, minimum width=3.4cm, minimum height=0.75cm, align=center, rounded corners=2pt, font=\scriptsize},
    safety/.style={rectangle, draw=red!80, fill=red!5, thick, minimum width=3.4cm, minimum height=0.75cm, align=center, rounded corners=2pt, font=\scriptsize},
    highlight/.style={rectangle, draw=teal!80, fill=teal!5, thick, minimum width=3.4cm, minimum height=0.75cm, align=center, rounded corners=2pt, font=\scriptsize},
    line/.style={draw, -latex, thick}
]
    \node[block] (state) {State Space $\mathcal{S}_j$\\Approach Queues $[q_l]_{l \in \mathcal{L}_j^{\text{in}}}$};
    \node[block, below=of state] (dqn) {SMDP Policy Engine\\Phase Selection $a_j \in \Phi_j$};
    \node[safety, below=of dqn] (bounds) {Action-Duration Bounds\\$10\,\text{s} \le \tau \le 45\,\text{s}, \; t_{\text{yellow}}=3\,\text{s}$};
    \node[block, below=of bounds] (actuation) {Physical Signal Actuation\\Phase Execution via TraCI};
    
    \node[highlight, right=1.0cm of dqn] (evp) {Rolling EVP Controller\\Preempt \& Momentum Flush};

    \path[line] (state) -- (dqn);
    \path[line] (dqn) -- (bounds);
    \path[line] (bounds) -- (actuation);
    \path[line] (state) -| (evp);
    \path[line, dashed, draw=teal!90] (evp) -- node[above, font=\tiny] {Priority} (dqn);
\end{tikzpicture}
\caption{The modular AURA control loop. The SMDP policy selects non-conflicting phase targets, action-duration bounds enforce clearance and stability invariants, and rolling preemption coordinates emergency transit with cross-street momentum recovery.}
\label{fig:architecture}
\end{figure}

\subsection{Semi-Markov Decision Process (SMDP) Formulation}
Standard discrete-time MDP formulations assume uniform transition durations ($\Delta t = 1$). In realistic traffic signal regulation, however, executing a phase switch incurs variable physical clearance times, minimum green intervals, and vehicle headway delays. We therefore formulate decentralized signal regulation as a **Semi-Markov Decision Process (SMDP)** defined by the tuple $\langle \mathcal{S}, \mathcal{A}, \mathcal{P}, \mathcal{R}, \mathcal{T}, \gamma \rangle$:

\subsubsection{Localized State Space ($\mathcal{S}$)}
For each signalized intersection $j$, the localized state vector $s_j(t) \in \mathcal{S}_j$ captures the vehicle halt count across its incoming approach lanes $\mathcal{L}_j^{\text{in}}$, the current active signal phase index $\phi_j$, and the elapsed phase duration $\tau_{\text{elapsed}}$:
\begin{equation}
s_j(t) = \Big( [q_l(t)]_{l \in \mathcal{L}_j^{\text{in}}}, \, \phi_j(t), \, \tau_{\text{elapsed}} \Big)
\label{eq:state}
\end{equation}
where $q_l(t)$ denotes the number of halting vehicles queued on incoming lane $l$. By defining the state over the topological lane set $\mathcal{L}_j^{\text{in}}$ rather than an assumed 4-way vector, AURA generalizes naturally to heterogeneous 3-way, 4-way, and asymmetric multi-lane urban intersections.

\subsubsection{Action Space ($\mathcal{A}$)}
The action space $\mathcal{A}_j = \Phi_j$ corresponds to the finite set of valid non-conflicting green phase allocations defined for intersection $j$:
\begin{equation}
a_j(t) \in \Phi_j
\label{eq:action}
\end{equation}
where each action chooses the next green phase target.

\subsubsection{Action-Duration Bounds and Transition Dynamics}
To eliminate phase flickering, guarantee pedestrian clearance, and prevent arterial starvation, action durations are bounded by a deterministic duration function $\tau(s_j, a_j)$:
\begin{equation}
\tau(s_j, a_j) = \max\Big(t_{\min}, \, \min\big(2 \cdot \max_{l \in \mathcal{L}_{a_j}^{\text{in}}} q_l(t), \, t_{\max}\big)\Big)
\label{eq:duration}
\end{equation}
where $t_{\min} = 10\,\text{s}$ enforces the statutory minimum clearance headway for queued vehicles to accelerate from rest, and $t_{\max} = 45\,\text{s}$ bounds maximum green duration against opposing arterial starvation. Whenever a phase transition occurs ($a_j \ne \phi_j$), a mandatory 3-second yellow clearance interval $t_{\text{yellow}} = 3\,\text{s}$ is deterministically executed prior to displaying the new green phase.

\subsubsection{Duration-Discounted Bellman Optimization}
Because the duration $\tau(s, a)$ varies per decision step, the standard fixed-step Bellman equation is mathematically invalid. We apply continuous-time discounting within the SMDP Q-learning update:
\begin{equation}
\mathcal{L}(\theta) = \mathbb{E}_{(s, a, r, \tau, s') \sim \mathcal{D}} \left[ \left( r + \gamma^{\tau} \max_{a' \in \Phi_j} Q(s', a'; \theta^-) - Q(s, a; \theta) \right)^2 \right]
\label{eq:smdp_loss}
\end{equation}
where $\gamma = 0.95$ denotes the base per-second discount factor, properly weighted by $\gamma^\tau$ across the variable duration $\tau \in [10, 45]$\,s, and $\theta^-$ represents a target network periodically synchronized every $C = 1{,}000$ steps. Transitions $(s, a, r, \tau, s')$ are stored in an experience replay buffer $\mathcal{D}$ of capacity $N = 10{,}000$ and sampled with mini-batch size $B = 64$.

\subsection{Reward Function and Post-Preemption Momentum Recovery}
The primary objective is maximizing network vehicle throughput while minimizing queue length and preventing cross-corridor gridlock. The scalar reward $r(s, a)$ comprises three terms:
\begin{equation}
r(s, a) = - \sum_{l \in \mathcal{L}_j^{\text{in}}} q_l(t) + \lambda \cdot \max\Big(0, \sum_{l \in \mathcal{L}_j^{\text{in}}} \big(q_l(t-\tau) - q_l(t)\big) \Big) + r_{\text{momentum}}(s)
\label{eq:reward}
\end{equation}
where the first term penalizes instantaneous halting vehicles, the second term provides a momentum throughput bonus ($\lambda = 1.5$) whenever queues decrease over the phase interval, and $r_{\text{momentum}}(s)$ represents an explicit shockwave mitigation term:
\begin{equation}
r_{\text{momentum}}(s) = - \beta \sum_{l \in \mathcal{L}_{\text{cross}}} q_l(t) \cdot \mathbb{I}(t_{\text{wait}, l} > t_{\text{thresh}})
\label{eq:momentum_term}
\end{equation}
where $\beta = 2.0$ heavily penalizes lingering queues on cross-street approaches following an emergency vehicle preemption event.

\subsection{Rolling Emergency Preemption and Corridor Coordination}
Unlike naive centralized preemption that forces entire multi-kilometer grids into gridlock, AURA implements a \textbf{Rolling Localized Preemption Protocol}:
\begin{enumerate}
    \item \textbf{Localized Proximity Trigger:} An intersection $j$ only preempts when an approaching emergency vehicle is detected physically on its incoming approach ($d \le 120\,\text{m}$), activating Phase $\phi_{\text{EV}}$ with guaranteed yellow clearance.
    \item \textbf{Sequential Green Wave:} Downstream signals coordinate sequentially as the EV traverses each intersection, maintaining continuous headway without paralyzing non-adjacent corridors.
    \item \textbf{Post-Preemption Momentum Flushing:} As soon as the EV exits intersection $j$, intersection $j$ terminates preemption and immediately activates Eq.~\eqref{eq:momentum_term}, prioritizing the starved cross-street queues for 25 seconds before returning to standard SMDP regulation. This dynamic flushing actively dissipates the secondary congestion shockwave.
\end{enumerate}
